\documentclass[../../../include/open-logic-section]{subfiles}

\begin{document}

\olfileid{sth}{story}{blv}

\olsection{Lampiran: Hukum Dhasar V Frege}

Ing \olref[rus]{sec}, kita nerangake yen Russell ngrumusake
paradokse minangka masalah tumrap sistem kang dijlentrehake
Frege ing \emph{Grundgesetze}. Sistem Frege ora ngemot rumusan
langsung Komprehensi Naif. Mula ing lampiran iki kita bakal
nerangake kanthi cekak banget apa kang \emph{pancen} ana
ing sistem Frege, sarta sesambungane karo Komprehensi Naif
lan Paradoks Russell.

Sistem Frege iku \emph{tataran kapindho} lan dirancang kanggo
ngrumusake gagasan \emph{ekstensi sawijining konsep}.\footnote{Yen
diandharake kanthi luwih trep, Frege nyoba ngrumusake gagasan
kang luwih umum: ``rentang nilai'' sawijining fungsi. Ekstensi
konsep iku kasus khusus saka gagasan kang luwih umum kasebut.
Kanggo rinci, delengen \citet[pp.\ 8--9]{Heck2012}.}
Kanthi notasi kang kailhami Frege, kita nulis
$\fregeext{x}{F(x)}$ kanggo \emph{ekstensi konsep~$F$}.
Iki piranti kang njupuk \emph{predikat}, ``$F$'', lan
ngowahi dadi \emph{term} (tataran kapisan), yaiku
``$\fregeext{x}{F(x)}$''. Kanthi piranti iki, Frege menehi
\emph{definisi} keanggotaan kaya mangkene:
\[
	a \in b =_\text{df} \exists G(b = \fregeext{x}{G(x)} \land Ga)
\]
kurang luwih: $a \in b$ yen lan mung yen $a$ kalebu ing konsep
kang ekstensine $b$. (Gatekna kuantor ``$\exists G$'' iku
tataran kapindho.) Frege uga nampa prinsip ing ngisor iki,
kang diarani \emph{Hukum Dhasar V}:
$$\fregeext{x}{F(x)} = \fregeext{x}{G(x)} \liff \forall x (Fx \liff Gx)$$
kurang luwih: konsep-konsep nduweni ekstensi kang identik yen
lan mung yen konsep-konsep iku koekstensif. (Maneh,
``$F$'' lan ``$G$'' loro-lorone ana ing posisi predikat.)
Saiki prinsip prasaja nyambungake keanggotaan karo pemenuhan sipat:

\begin{lem}[ing \emph{Grundgesetze}]\ollabel{lem:Fregeextensions}
$\forall F \forall a(a \in \fregeext{x}{F(x)} \liff Fa)$
\end{lem}

\begin{proof} 
Tetepna $F$ lan $a$. Saiki $a \in \fregeext{x}{F(x)}$ yen lan
mung yen $\exists G(\fregeext{x}{F(x)}
= \fregeext{x}{G(x)} \land Ga)$ (miturut definisi keanggotaan),
yen lan mung yen $\exists G(\forall x(Fx \liff Gx) \land Ga)$
(miturut Hukum Dhasar V), yen lan mung yen $Fa$ (miturut
logika tataran kapindho kang dhasar).
\end{proof}

Lan iki meh langsung ngasilake Komprehensi Naif:

\begin{lem}[ing \emph{Grundgesetze}.]
$\forall F \exists s \forall a (a \in s \liff Fa)$
\end{lem}

\begin{proof}
Tetepna $F$; miturut \olref{lem:Fregeextensions} kita nduweni
$\forall a (a \in \fregeext{x}{F(x)} \liff Fa)$; mula
$\exists s\forall a(a \in s \liff Fa)$ kanthi generalisasi
eksistensial. Asil iki lumaku amarga $F$ dipilih sembarang.
\end{proof}

Paradoks Russell tuwuh yen $F$ dipilih kanthi
$\forall x(Fx \liff x \notin x)$.

\end{document}
